% D21
% Auto-generated from markdown source
% origin: ai

\documentclass[12pt, a4paper]{article}
\usepackage[utf8]{inputenc}
\usepackage[T1]{fontenc}
\usepackage{amsmath, amssymb, amsfonts, amsthm}
\usepackage{textcomp}
\usepackage{geometry}
\usepackage{hyperref}
\usepackage{booktabs}
\usepackage{graphicx}
\usepackage{fancyhdr}
\usepackage{tcolorbox}
\usepackage{enumitem}

\geometry{margin=1in}
\pagestyle{fancy}
\fancyhf{}
\rhead{D21}
\lhead{\thepage}

\newtheorem{definition}{Definition}[section]
\newtheorem{theorem}{Theorem}[section]
\newtheorem{corollary}{Corollary}[section]
\newtheorem{proposition}{Proposition}[section]
\newtheorem{lemma}{Lemma}[section]
\newtheorem{conjecture}{Conjecture}[section]

\newcommand{\calI}{\mathcal{I}}
\newcommand{\calP}{\mathcal{P}}
\newcommand{\calE}{\mathcal{E}}
\newcommand{\calS}{\mathcal{S}}
\newcommand{\calH}{\mathcal{H}}
\newcommand{\calO}{\mathcal{O}}
\newcommand{\calF}{\mathcal{F}}
\newcommand{\calU}{\mathcal{U}}
\newcommand{\calV}{\mathcal{V}}
\newcommand{\calA}{\mathcal{A}}
\newcommand{\calG}{\mathcal{G}}
\newcommand{\calB}{\mathcal{B}}
\newcommand{\calC}{\mathcal{C}}
\newcommand{\calL}{\mathcal{L}}
\newcommand{\calM}{\mathcal{M}}
\newcommand{\calX}{\mathcal{X}}
\newcommand{\calY}{\mathcal{Y}}
\newcommand{\calW}{\mathcal{W}}
\newcommand{\calT}{\mathcal{T}}
\newcommand{\calZ}{\mathcal{Z}}
\newcommand{\calR}{\mathcal{R}}
\newcommand{\calQ}{\mathcal{Q}}
\newcommand{\calD}{\mathcal{D}}
\newcommand{\calN}{\mathcal{N}}
\newcommand{\calK}{\mathcal{K}}

\newcommand{\Tr}{\text{Tr}}
\newcommand{\E}{\mathbb{E}}
\newcommand{\diff}{\mathrm{d}}
\newcommand{\KL}{D_{\mathrm{KL}}}
\newcommand{\argmax}{\arg\max}
\newcommand{\argmin}{\arg\min}
\newcommand{\hbarF}{\hbar}

\newtcolorbox{thmbox}[1]{
  colback=blue!5!white, colframe=blue!75!black, fonttitle=\bfseries,
  title={#1}, breakable
}
\newtcolorbox{defbox}[1]{
  colback=green!5!white, colframe=green!75!black, fonttitle=\bfseries,
  title={#1}, breakable
}

\begin{document}

\title{D21}
\author{RALPH Loop}
\date{}
\maketitle

\begin{abstract}
This paper presents D21, formalizing the underlying principles using information-theoretic methods.

\textbf{Keywords}: D21, information theory, formalization.
\end{abstract}

\tableofcontents
\section{D21}




\subsection{0.}


homology groups Betti holes\textbf{*}***

\textbf{D21 }**** Betti
\begin{itemize}
\item ****
\item ****
\item ****
\end{itemize}




\subsection{1.}


\subsubsection{1.1}


\textbf{ D21-1 \(\mathcal{M}\)}
\(\mathcal{S}\)****
\[\mathcal{M} = \bigcup_{x \in \mathcal{S}} B_\epsilon(x)\]

\(B_\epsilon(x)\)  \(\epsilon\)


\subsubsection{1.2 Simplicial Complex  Simplicial Homology}


\textbf{ D21-2 \(K\)}
\(\mathcal{M}\) \textbf{\(\epsilon\)-}
\[K_\epsilon = \{[v_0, \ldots, v_k] \mid d(v_i, v_j) \leq \epsilon, \forall i, j\}\]

\([v_0, \ldots, v_k]\)  \(k\)- \(k\)-

****
\[H_k(K) = \ker(\partial_k) / \text{im}(\partial_{k+1})\]

\(\partial_k\)


\subsubsection{1.3 Betti}


\textbf{ D21-3Betti  \(\beta_k\)}
\[\beta_k = \text{rank}(H_k(K)) = \dim(\ker(\partial_k) / \text{im}(\partial_{k+1}))\]

****raw.md \#55
| Betti  |  |  |
|----------|----------|----------|
| \(\beta_0\) |  |  |
| \(\beta_1\) |  |  |
| \(\beta_2\) |  |  |
| \(\beta_k\) | \(k\)- | \(k\)- |

****
\[\beta_k = b_k - a_{k+1}\]
\(b_k = \dim(\ker \partial_k)\) \(a_{k+1} = \dim(\text{im } \partial_{k+1})\)


\subsection{2.}


\subsubsection{ D21-1 Betti}


****
**** \(\mathcal{M}\)  \(\epsilon\)- \(K_\epsilon\)

\textbf{D21-1 }
\[\beta_0(K_\epsilon) = 1 \iff \mathcal{M} \text{ }\]

\textbf{ D21-1a}
\begin{itemize}
\item \textbf{\(\beta_0 = 1\)}
\item \textbf{\(\beta_0 > 1\)}****
\end{itemize}


\textbf{D1-S12 }
\[\beta_0 = M \iff \text{chimera  } M \text{ }\]


\subsubsection{ D21-2 =}


****
\textbf{*}*** 1-
\[\gamma = \sum_{i=1}^{n} [v_i, v_{i+1}], \quad \gamma \in Z_1 \setminus B_1\]

\textbf{D21-2 }
\[\beta_1(K_\epsilon) > 0 \implies \exists \gamma: \gamma \neq \partial \sigma, \quad \forall \sigma \in C_2\]

\textbf{ D21-2a} D1-S12
\[\beta_1 > 0 \iff \text{}\]

****
\begin{itemize}
\item \textbf{ = 0}
\item \textbf{ > 0}****
\end{itemize}



\subsubsection{ D21-3}


****
**** 2- 2-
\[\sigma \in C_2, \quad \partial_2(\sigma) = 0, \quad \nexists \tau \in C_3: \partial_3(\tau) = \sigma\]

\textbf{D21-3 }
\[\beta_2(K_\epsilon) > 0 \implies \exists \text{ } \Omega_{\text{blind}}: \partial \Omega_{\text{blind}} \in \mathcal{M}, \quad \Omega_{\text{blind}} \cap \mathcal{M} = \emptyset\]

\textbf{ D21-3a}
\begin{itemize}
\item \textbf{ = 0} 2-
\item \textbf{ > 0}****
\end{itemize}


\textbf{**  }**


\subsubsection{ D21-4}


****
\textbf{*}*** \(K_\epsilon(t)\)

\textbf{D21-4 }
\[\frac{d\beta_k}{dt} = \begin{cases} \geq 0 & \text{}\\ \leq 0 & \text{} \end{cases}\]

****
\[t_c = \inf\{t \mid \beta_k(t) \neq \beta_k(t - \delta t)\}\]

\textbf{ D21-4a} D1-S12
\textbf{\(t_c\) }S12-1 ****


\subsubsection{ D21-5Betti}


****
\textbf{** \(\epsilon \in [0, \infty)\) }Vietoris-Rips **
\[K_\epsilon \subseteq K_{\epsilon'} \quad \forall \epsilon < \epsilon'\]

****
\[H_k^{persistence} = \{(t_{\text{birth}}, t_{\text{death}}) \mid \beta_k(t_{\text{birth}}) > \beta_k(t_{\text{death}})\}\]

\textbf{D21-5 }
\[\exists \Delta > 0: \forall \epsilon: \beta_k(K_\epsilon) = \beta_k(K_{\epsilon+\Delta}) \iff \mathcal{M} \text{  } k\text{-}\]

\textbf{persistence diagram}
\[PD_k = \{(t_i^{b}, t_i^{d})\}_{i=1}^{n} \subset [0, \infty)^2\]

\textbf{ D21-5a}
\begin{itemize}
\item ****\(t_d - t_b\) =
\item **** =
\end{itemize}



\subsection{3.}


\subsubsection{ D21-A}


\textbf{ D21-1  D21-2}
1.  LLM cosine similarity
2.  Vietoris-Rips \(\epsilon\)-
3.  Betti  \(\beta_0, \beta_1, \beta_2\)
4. **** \(\beta_1 > 0\)


\subsubsection{ D21-BBetti}


\textbf{ D21-4}
1.
2.  Betti
3. ****


\subsubsection{ D21-C}


\textbf{ D21-5}
1.
2.  Betti
3. ****


\subsubsection{ D21-D-}


\textbf{ D21-3}
1.
2.   > 0
3. ****   0


\subsection{4.}


\subsubsection{4.1 D1-S12}


\textbf{D21-D1-S12 }
\begin{itemize}
\item D1-S12  chimera M=2  D21  \textbf{  1}
\item D1-S12  N   D21  ****
\item \textbf{** = } > 0**
\end{itemize}



\subsubsection{4.2 D9-IV}


\textbf{D21-D9-IV }
\begin{itemize}
\item D9-IV \textbf{ T}  D21  \textbf{ }
\item D9-IV \textbf{**  D21  } **
\item \textbf{** = }**
\end{itemize}



\subsubsection{4.3 D16}


\textbf{D21-D16 }
\begin{itemize}
\item D16 \textbf{**  D21  }Vietoris-Rips  **
\item D16 \textbf{**  D21  }**
\item \textbf{** = }**
\end{itemize}



\subsection{5.}


\subsubsection{5.1}


1. \textbf{D21-1 }  > 1
2. \textbf{D21-2 } > 0
3. \textbf{D21-3 }
4. \textbf{D21-5 }

\subsubsection{5.2}


1. **** N > 10^4
2. ****Betti
3. ****


\subsection{6.}


\#\#\#

**** \(\mathcal{M}\)  \(\epsilon\)-Vietoris-Rips  \(K_\epsilon\) \(H_k(K_\epsilon)\)

****
1. ****\(\beta_0 = \text{rank}(H_0) = \#(\text{})\)
2. ****\(\beta_1 = \text{rank}(H_1) \geq 1 \implies \text{}\)
3. ****\(\beta_2 = \text{rank}(H_2) > 0 \implies \text{}\)
4. ****\(t_c = \inf\{t \mid \beta_k(t) \neq \beta_k(t-\delta t)\}\)
5. ****\(\beta_k(K_\epsilon) = \beta_k(K_{\epsilon+\Delta}) \iff \text{}\)

\textbf{*}\textbf{*}\textbf{Betti }**


*D21  Betti *

\begin{thebibliography}{99}
\bibitem{ref1} Shannon, C.E. (1948). A Mathematical Theory of Communication.
\bibitem{ref2} Cover, T.M. (2006). Elements of Information Theory.
\end{thebibliography}

\end{document}